\section{Sections moved from the main paper}
\label{sec:moved}
The two transfer checks below were part of the first version of the main paper and are kept here unchanged.

\begin{figure}[t]
\centering
\includegraphics[width=\linewidth]{branches.pdf}
\caption{Which SALAD branch separates a fault from a valid edit? (a)~Juice edit: a same-type (valid) or other-type (fault) icon, pasted identically. (b,c)~Scenes alarmed per branch at its own 10\% normal cutoff (dashed); on pushpin composites, global RGB and fusion also alarm on most untouched sources. (d)~Juice fault-minus-sham alarm gap, 95\% intervals. Photos: MVTec LOCO AD, CC BY-NC-SA 4.0.}
\label{fig:branches}
\end{figure}
\subsection{Branch selection under two rules}
\label{sec:transfer}
SALAD fuses three branches: a local RGB appearance map, the composition map studied above, and a global RGB score that pools student features over the scene and over regions defined by the part map~\cite{fucka2025}. Our replay of the complete released system reproduces the published 100.0\% logical AUROC on 91 native logical pushpin faults and 138 normal photographs.

On the matched pushpins, however, the shams reveal a cost that this AUROC does not capture (\cref{fig:branches}b). Fusion ranks every fault above its sham but raises an alarm on 130 of 138 extracted shams and 75 of 138 sources, with 55 shams alarming although their source does not. The local RGB branch alone alarms on 128 shams and only 9 sources; since both arms contain the same compositing, it responds to the edit rather than to the rule. The composition branch alarms on 7 shams, and an explicit one-pin-per-compartment rule applied to the same extracted maps raises no sham alarm while detecting every fault.

A second rule shows that the informative branch depends on the rule. MVTec LOCO requires the icon on the upper label of a juice bottle to match its liquid~\cite{bergmann2022}. We retain 89 of 94 eligible normal test bottles and paste intact icons from a disjoint bank of training photographs into the same label region (\cref{fig:branches}a); an icon of the same type yields a valid sham and an icon of another type a fault. Both use the same offset, blending and footprint, and their mean blended RGB changes differ by at most half a grey level. Since the liquid, the bottle and the part map are unchanged, the composition branch returns identical maps for all three variants. The global RGB branch ranks all 89 faults above their shams and alarms on every fault but only 14 shams, a fault-minus-sham alarm gap of $+84.3$ points [76.4, 91.0] (\cref{fig:branches}d), slightly above that of fusion ($+79.8$ [70.8, 87.6]). The local RGB branch ranks the upper label first among five bottle regions for 74.2\% of faults but also for 85.4\% of shams, and alarms on 25 faults and 28 shams. The two rules therefore favour different branches: composition evidence decides placement in a tray, whereas global appearance captures a label--liquid correspondence that the part map cannot express.

\begin{figure}[t]
\centering
\includegraphics[width=\linewidth]{graph_audit.pdf}
\caption{A masked GNN responds to switched links through node degree (78 GlaS fields, 12 donors). (a)~Edited-graph AUROC; local spacing adds 15.3 points to typed distance. (b)~AUROC of the change between matched graphs; hollow: 50 by construction. (c)~Top-four recovery of the masked GNN as fitted, with degree inputs fixed, and refit with mean aggregation. Donor-paired 95\% intervals.}
\label{fig:graph}
\end{figure}
\subsection{Response evidence in part graphs}
\label{sec:graph}
The matched design also applies to graphs of parts. We build gland graphs from GlaS histology images~\cite{glas2017} and cardiac structure graphs from HVSMR-2.0~\cite{pace2024}, and apply \emph{degree-preserving switches} that remove two links and add two others so that every node keeps its degree. Readouts are evaluated by how well they rank the four switched pairs. On GlaS, adding local gland spacing to typed centroid distance improves the edited-graph AUROC by 15.3 points [10.9, 20.1], and the gain persists for additional length-matched switches and under joint matching of raw and local distances ($+16.8$ [11.3, 23.0]); these tiers reuse donors and involve no pathological labels.

For a candidate link $e=(u,v)$ between nodes $u$ and $v$, let $z_e$ be the predicted logit and $a_e\in\{0,1\}$ the observed link state. For a link predictor, the Bernoulli loss $\operatorname{softplus}(z_e)-a_ez_e$ changes with $a_e$ even when the evidence $z_e$ is fixed. A pair MLP observes only geometry and roles, so its probability cannot change and it ranks by response at chance level (\cref{fig:graph}b). A query-masked neighbourhood model has an additional route: removing the queried link before computing the degree $d_u$ of an endpoint yields $d_u-a_e$, so a switch changes the masked degree at the four target pairs although both graphs share the same degree sequence, an effect related to the target-link exclusion studied in~\cite{zhu2024spot}. Ranked by absolute change in probability, the masked GNN recovers 87.5\% of the switched pairs among its top four (\cref{fig:graph}c). Holding its degree inputs fixed reduces this to 33.7\%, and refitting with mean neighbourhood aggregation and no degree inputs on the same four GlaS folds, which reaches 95.7--98.2\% training AUROC on normal links, reduces it to 46.6\%, 40.9 points lower [32.8, 49.8]. Since mean aggregation still admits indirect cues such as an empty neighbourhood, both interventions indicate a degree pathway, and recovery of known graph edits should not be interpreted as evidence of validity.

